\section{Agentic EQA tools}
\label{app:agentic_tools}

This appendix documents the scored HM-EQA / Dynagraph \emph{agentic} tool set,
distinct from interactive CHAT. Canonical names, descriptions and argument schemas
live in \texttt{emet.agent.skills.specs}; EQA bindings live in
\texttt{emet.memory.graph\_eqa.agentic.tools}. The older
\texttt{agentic\_tools} import remains a compatibility shim.
When \texttt{eqa.agentic\_verify} (or \texttt{EMET\_EQA\_AGENTIC\_VERIFY=1}) is set, \texttt{GraphEQAController.run\_eqa} runs \texttt{AgenticEQAExecutor} instead of the classic planning/\texttt{query\_answer} loop (Section~\ref{sec:eqa_loops}).

The historical scored rows pin their own source/configuration; current state,
progress and recovery mechanisms below must not be assumed enabled in every
archived experiment. Appendix~\ref{app:combined_agent} summarizes the separate
CHAT/TAMP orchestrator and its capability gates.

\subsection{Protocol}

\paragraph{System prompt.}
A fixed identity + tool catalog + JSON response-format block (byte-stable across rounds for prefix KV caching).
The model must reply with only:
\begin{verbatim}
{"tool_calls": [{"name": "<tool>", "arguments": {...}}, ...], "message": ""}
\end{verbatim}
The actual prompt also requests \texttt{current\_room} and
\texttt{in\_target\_area} context. In \texttt{model\_final}, the model selects
the next observation/action from images, history and measured outcomes, without
a required investigate/explore sequence. Explicit submitted semantic answers are
preserved; omitting the answer requests a final model evidence read. Historical
\texttt{legacy} and \texttt{grounded\_v2} policies retain their own answer gates,
progress rules and action substitutions. None of these semantic decisions alone
admits object geometry or authorizes motion/manipulation.
SigLIP/OWL are intentionally high-recall (high false positive) support for choosing where to look---not a precision oracle (Appendix~\ref{app:agentic_verify_gate}).

\paragraph{State message (Investigate / Explore cards).}
Each round appends question, budgets, graph counts, and split card lists: \textbf{Investigate} (graph / confirmed / SigLIP) with per-query \texttt{investigated=$N$}, closest approach, and recent station outcomes; \textbf{Explore} (frontiers, coverage only). Grounded state schema v3 renders recent work as normalized intent, stable place/frontier/view identity, concrete approach/view variant, typed outcome, and material-progress reason; the mutable observation number is retained only as a trailing tool adapter.
Retrieval only decides which cards fit in context ($K{=}6$ by default via \texttt{EMET\_EQA\_HYP\_RECALL\_K}); the router is not given a composite navigation-utility score.
Visited frontier nodes are retired from the graph so they are not re-offered as frontiers.

\paragraph{Automatic capture.}
After motion tools succeed, the executor captures a fresh observation and updates voxel + graph memory without exposing a separate \texttt{capture\_and\_update} tool to the VLM.

\paragraph{Investigate vs.\ explore.}
The router chooses explicitly: \texttt{investigate(obs\_id)} for a closer look at a listed place card (graph / confirmed / SigLIP), or \texttt{explore\_frontier} to grow coverage when place approaches are exhausted or no place is worth pursuing.
State splits \textbf{Investigate} and \textbf{Explore} cards; Investigate lines show per-query \texttt{investigated=$N$}, \texttt{approaches=$k$/$N$}, \texttt{coverage=open|closed} (object footprint $\cap$ unexplored frontier), closest approach, and recent station verify/assess outcomes.
\texttt{navigate\_to\_obs} remains a callable alias of \texttt{investigate}, but is omitted from \texttt{model\_final} advertisement.
Optional \texttt{toward} on explore is coverage bias only---not a substitute for investigating a place card.

\paragraph{Station observations are not place cards.}
Automatic capture stations (\texttt{\_station\_obs\_ids}) appear in memory but are \emph{excluded} from Investigate cards and blocked as \texttt{investigate} targets (\texttt{STATION\_OBS\_NOT\_PLACE}).
Chasing station obs along a patio (e.g.\ q105) burned the nav budget with zero \texttt{explore\_frontier} calls; the paper-router stack filters them so the router must pick real graph / SigLIP place cards or explore.

\paragraph{Historical explore-after-ABSENT policy.}
Outside \texttt{model\_final}, after a close approach that returns detector/VLM \texttt{ABSENT}, the executor sets \texttt{\_prefer\_explore} and forces one \texttt{explore\_frontier} (state nudge + fallback + run-loop redirect) instead of immediately re-investigating another nearby card.
When action history is agent-visible, state also surfaces a short \textbf{Recent actions} history so the router sees prior ABSENT / explore outcomes within the episode.

\paragraph{Configurable action-history visibility.}
For \texttt{grounded\_v2}, \texttt{eqa.attempt\_ledger\_mode} is a three-way,
manifest-frozen control: \texttt{off} neither persists dedicated attempt rows nor
renders action history; \texttt{shadow} persists the same rows and artifacts but
hides them from the router; and \texttt{agent} exposes recent action outcomes,
navigation-loop flags, per-place attempt summaries/failure risk, and global
attempt rows. Room history and live approach affordances remain separate state
axes. Strict complete variants are checked in as
\texttt{configs/benchmarks/hmeqa\_action\_history\_shadow.yaml} and
\texttt{hmeqa\_action\_history\_agent.yaml}; the run manifest records their
effective values and SHA-256 source.

\paragraph{Static-world semantic progress policy.}
Action-history visibility is independent from the manifest-frozen
\texttt{eqa.action\_progress\_mode} axis. \texttt{off} preserves the legacy
candidate policy; \texttt{shadow} computes and traces counterfactual
suppression decisions without changing cards or execution; and
\texttt{enforce} omits an equivalent terminal/no-progress candidate before
routing while its action-specific target state is unchanged. Equivalence uses
stable semantic work plus the concrete investigate approach, immutable verify
view and phrase, or material frontier geometry/goal. New approaches, views,
target-local evidence, coverage/geometry, and partial navigation progress
reopen work. Checked-in controls are
\texttt{hmeqa\_action\_progress\_shadow.yaml} and
\texttt{hmeqa\_action\_progress\_enforce.yaml}.

This policy is not a permanent failed-action blacklist. Its first version is
deliberately scoped to HM-EQA's mostly static scenes. In a dynamic embodied
setting, target/environment change events, map revision, elapsed-time/TTL, and
evidence staleness must invalidate or decay a suppression; a future scheduler
may cool down or re-rank actions instead of removing them. We therefore keep
the default \texttt{off} and treat the shadow/enforce comparison as a
mechanism/efficiency diagnostic, not evidence for a general lifelong-memory
policy.

\paragraph{Earlier action-history visibility diagnostic (2026-08-23).}
To isolate the visibility channel from the earlier bundled grounded-state
comparison, we use a paired four-question diagnostic (\{6,11,12,47\}) in which both arms use
\texttt{grounded\_v2} with agent-visible graph and room evidence, no investigate
stamp, no HM3D semantic sensor, and no per-question enriched labels. Only
\texttt{attempt\_ledger\_mode=shadow} versus \texttt{agent} changes. Eight
episodes are still too small for an accuracy claim; the test is reported as
router-state/process evidence.

At clean commit \texttt{63abd929}, both managed runs passed detached EGL probes,
completed without native crashes, and published validated four-unit completion
records. Shadow and agent-visible history both scored \textbf{3/4}, with the
same per-question letters (q6 B, q11 A, q12 D, q47 A). Mean planning steps
increased from 35.25 (shadow) to 39.25 (agent-visible), with one budget hit in
each arm. The predefined ledger repeat-failure count fell from $6$ to $4$
(q11: $5\!\rightarrow\!3$), but total ledger attempts rose from $31$ to $37$
and q12 required two extra decision rounds / 16 extra planning steps. All 12
shadow router states kept the dedicated history sections empty; recent actions,
persisted loop flags, and global attempts were non-empty in 10/14 treatment
router calls. Thus the visibility mechanism and shadow no-leakage control
worked, but this small diagnostic does not support default enablement or scale.
Frozen details:
\texttt{paper/data/hmeqa\_agentic\_h2h/action\_history\_pair\_20260823.json}.
A separate matched action-progress diagnostic on the same four questions
(\texttt{shadow} versus \texttt{enforce}, ledger visibility held at
\texttt{agent}) completed native-clean with schema-v4 manifests and gate traces.
Shadow preserved behavior and matched the prior $3/4$ baseline (mean planning
steps $35.25$); \texttt{enforce} regressed to $2/4$ with higher mean steps
($43.25$) and did not fire explicit suppress dispositions on this slice. Keep
\texttt{action\_progress\_mode} off by default; use \texttt{shadow} for
diagnostics only. Frozen details:
\texttt{paper/data/hmeqa\_agentic\_h2h/action\_progress\_pair\_20260823.json}.

\subsection{Answer-mode tools}

\begin{description}
  \item[\texttt{read\_history}] Read exact retained text by archived
    \texttt{message\_id}, with optional \texttt{offset}/\texttt{max\_characters}
    pagination. This does not acquire new evidence.
  \item[\texttt{recall\_image}] Reattach retained historical pixels by
    \texttt{image\_id}; recall is not fresh visibility or current geometry.
  \item[\texttt{inspect\_graph}] Refresh question keywords and retrieve up to $K$ place evidence cards from the scene graph, SigLIP hits, and frontier waypoints (diversified source packing).
    Use when the evidence list looks stale or empty.
    Parameters: none.
  \item[\texttt{explore\_frontier}] Navigate for new coverage; legacy policies can prefer this after exhausted Investigate cards or close ABSENT.
    Optional \texttt{toward} is weak coverage bias only---not investigate.
    Agentic prefers a VLM pick among reachable frontier RGBs when available.
    Successful visits retire nearby frontier graph nodes.
  \item[\texttt{investigate}] Closer look at a listed Investigate \texttt{obs\_id}: sample a navigable floor pose in an annulus around the object (prefer frontier-adjacent), look around, verify/assess; card shows \texttt{coverage=open|closed} from footprint$\cap$frontier and \texttt{approaches=$k$/$N$}.
  \item[\texttt{navigate\_to\_obs}] Compatibility alias for \texttt{investigate}, hidden in \texttt{model\_final}.
  \item[\texttt{look\_around}] Refresh evidence with the configured scan policy;
    when sweeping is disabled it captures the current heading only. Optional
    \texttt{pan\_rad}/\texttt{tilt\_rad}, advertised only with declared head
    limits, request a measured stationary head view instead of a sweep. Omitted
    joints retain their measured angle; out-of-range requests are rejected.
  \item[\texttt{verify\_siglip}] Cheap visual check: does string \texttt{phrase} match the current camera view or stored view \texttt{obs\_id} (default $-1$ = latest fresh view)?
    Reports retrieval proposals separately from VLM visual presence, answerability
    and object-geometry admission. Positive presence is not a geometric position;
    voxel detection handles are not camera-view IDs.
  \item[\texttt{submit\_answer}] Submit optional semantic \texttt{answer} text
    (for MCQ, option text without its letter). \texttt{model\_final} preserves an
    explicit submission; omission requests a final model read. Historical policies
    can require positive answer evidence or apply their configured unverified
    fallback. Detector PRESENT alone is not semantic proof. OVMM localization
    independently requires admitted object geometry.
\end{description}

\subsection{Explore-mode tools}

When no question is provided (map-building explore), \texttt{submit\_answer} is replaced by:

\begin{description}
  \item[\texttt{finish}] End exploration with an optional one-sentence \texttt{summary}.
    Completion follows the selected episode policy; this tool is mutually exclusive with \texttt{submit\_answer}.
\end{description}

\subsection{Capability-qualified shared tools}

EQA can additionally bind CHAT's \texttt{observe\_floor},
\texttt{rotate\_base} and \texttt{move\_forward} when the controller declares
support. Floor inspection requires measured head motion and calibrated depth,
rechecks the rejected footprint and restores task gaze; it does not authorize a
route. Relative motion reports actual progress and refusal reasons rather than
assuming requested travel. The registry contains ten EQA names across both
terminal submodes. The advertised \texttt{model\_final} pack contains eight base
tools, plus up to three shared tools; other policies retain the navigation alias.
The session schema is authoritative for optional arguments and bounds.

\subsection{Budgets and configuration}

Typical knobs (YAML \texttt{eqa.*} / env overrides): max agentic rounds, max nav steps per episode (\texttt{agentic\_max\_nav\_steps}), hyp recall $K$ (\texttt{EMET\_EQA\_HYP\_RECALL\_K}), SigLIP verify similarity floor (\texttt{verify\_min\_sim}), action-history visibility (\texttt{attempt\_ledger\_mode}), static-world retry policy (\texttt{action\_progress\_mode}), and trace collection (\texttt{EMET\_EQA\_TRACE}).
Habitat H2H entrypoint: \texttt{scripts/run\_hmeqa\_agentic\_h2h.sh} / \texttt{emet hmeqa h2h}.
Episode review: \texttt{emet hmeqa inspect OUT --qid N}.
